\section{Discussion}\label{discussion}

\subsection{Why auditable tree ensembles perform
strongly}\label{why-auditable-models-perform-strongly}

The central empirical finding is that auditable tree ensembles occupy the
leading observed point-estimate group and are competitive with deep-learning
alternatives under strict leave-one-condition-out evaluation.  In the
primary single-seed planned comparison, corrected statistical comparisons
established clear superiority only over the implementation-specific historical
shallow MLP baseline, while
most pairwise differences against deep-learning models were favourable by
point estimate but not significant.  A secondary current fixed-setting
five-seed sensitivity produces the same qualitative ordering for the selected
repeated baselines. These runs use the same canonical 14-condition
training budget for every repeated model, the current random forest is
significantly lower than the six repeated deep-learning baselines under
Holm--Bonferroni-corrected Wilcoxon tests on the 16 condition-level paired
MAEs (Supplementary Table~\ref{tab:supp-seed-repeat-all}). They do not exactly
reproduce the historical tuned neural checkpoints, so this sensitivity mixes
seed variation with protocol/configuration differences relative to the
historical artifacts. It is reported as current-protocol robustness evidence,
not as independent confirmation or the primary inferential conclusion. The observed RF advantage is most
plausibly related to dataset size, spectral feature structure, and the
grouped validation protocol.

\subsubsection{Small-data bias--variance
behaviour}\label{small-data-bias-variance}

The dataset contains only 319 samples spread across 16 grinding
conditions. With roughly 20 samples per condition, deep networks with
86--171k gradient-trained weights may have more capacity than the
319-sample dataset can reliably constrain. Random forests and LightGBM, by contrast,
partition the feature space with a large collection of shallow, axis-parallel
rules; they do not have gradient-trained neural weights and instead store tree counts, depths, leaves, split thresholds, feature indices, and leaf predictions. The observed ranking may reflect differences in inductive bias, regularisation, optimisation stability, and effective complexity under the small condition-level sample size:
Figure~\ref{fig:discussion-complexity} shows that tree ensembles achieve the lowest mean LOGO MAE, while the comparison uses checkpoint size as a storage-footprint proxy rather than a one-parameter stand-in for trees. ResNetVibCNN and the bilinear fusion network still require many trainable weights yet deliver higher error (0.027 and 0.032~\(\upmu\)m versus 0.020~\(\upmu\)m for the best random forest). This should be interpreted as a hypothesis
because architecture design and regularization could change the result; it
is consistent with broader evidence that tree-based models remain
competitive on small tabular or pre-processed spectral datasets
\cite{grinsztajn2022tree, shwartz2022tabular}.

\begin{figure}[htbp]
\centering
\includegraphics[width=140mm]{discussion_complexity_accuracy.png}
\caption{Model size versus LOGO accuracy. Tree ensembles (Random Forest, LightGBM) are compared with deep models using checkpoint size on a log scale rather than a one-parameter stand-in; lower MAE does not imply smaller checkpoints.}\label{fig:discussion-complexity}
\end{figure}

The fold-level variance seen in Table~\ref{tbl:statistical} reinforces
this interpretation. Even when the random forest has a lower point-estimate
MAE, the wide per-fold confidence intervals prevent most pairwise
differences from reaching significance. This variance likely reflects condition-to-condition heterogeneity,
potentially mixed with run-order or wheel-state effects; the present deep
models did not reduce this heterogeneity relative to the regularised tree
ensembles.

\subsubsection{Inductive bias of tree ensembles on spectral
features}\label{inductive-bias}

Spectrograms and scattering transforms are naturally organised as
frequency-indexed tables. A tree ensemble can place a split directly on a
single frequency bin, making it easy to isolate narrow bands such as the
20.9~kHz vibration bin identified by TreeSHAP in
Section~\ref{sec:top-features}; AE bins above 1~MHz should be treated as
empirical model attributions rather than as independently validated
physical features. The evaluated CNN architectures apply translation-invariant filters
and pooling operations that may blur some fine-grained frequency-bin
information. The sparse attribution profile suggests that narrow-bin
information may be useful, but controlled feature and architecture
ablations are needed to verify this.

In this dataset, per-sample dB-z and log-mel preprocessing appear to produce feature spaces in which tree splits are effective. Their canonical RF MAEs differ by only 0.0001~\(\upmu\)m, while the current documented dB-z pipeline does not exactly reproduce the 0.0200~\(\upmu\)m artifact. The evidence therefore supports a top group rather than a unique representation winner. In the implemented log-mel cache, local dB maps are averaged before inverse conversion to power and mel filtering. The recovered linear quantity is consequently a geometric-mean power representation, not the arithmetic mean obtained by mel-filtering each local power map before pass averaging. This non-standard processing order may suppress transient energy and should be considered when comparing the representation with conventional log-mel pipelines. Per-sample standardisation and mel filtering may alter the relative weighting of high-energy AE components and lower-energy vibration modes; controlled preprocessing ablations would be needed to verify this mechanism.
Wavelet scattering achieves a similar stabilisation through
mathematically-defined translation and deformation invariants. However, because our WST features (\texttt{ae\_wst}, \texttt{vib\_wst}) are computed by applying the 2D spatial scattering transform on top of the already-computed 2D STFT spectrogram (rather than directly on the raw 1D time-domain signals), there is a significant mathematical redundancy. The STFT magnitude operation has already discarded phase information and performed local temporal averaging, meaning that the subsequent scattering transform acts as a secondary layer of modulation filtering rather than a raw signal representation. This redundancy, combined with boundary smoothing on the short vibration time series ($T=13$), may contribute to the observed WST ranking, but controlled ablations would be needed to establish the mechanism. In this dataset, simple learners on engineered spectral
representations achieve lower point-estimate errors than the evaluated
neural baselines---a pattern that has been observed in other small-data
spectral tasks and is reproduced here for precision grinding.

\subsubsection{LOGO as a proxy for operational generalisation and the
risk of CV leakage}\label{logo-operational-generalisation}

Many prior studies in tool-condition and surface-quality monitoring use
ordinary K-fold cross-validation, which can place samples from the same
grinding condition in both training and test folds. Because a condition
shares wheel state, coolant history, dressing state, and perhaps thermal
history, this protocol leaks condition-specific artefacts across folds and
overestimates real-world performance.

LOGO cross-validation treats each condition as a unit, so the test fold
always contains a previously unseen grinding state. This is a stricter proxy
for transfer to an unseen process-parameter combination within the same machine, wheel campaign, workpiece material, and sensor installation. The auditable tree ensembles achieved lower point-estimate errors under this protocol, possibly
because they did not overfit to the idiosyncrasies of individual
conditions. The persistent influence of
Condition~7 (Section~\ref{sec:conditions}) is precisely the kind of
condition-level heterogeneity that LOGO exposes and that ordinary K-fold
would dilute.

A final caveat concerns model selection. The deep-learning entries were
optimised by random search over model and training hyperparameters on the
same outer LOGO folds, whereas the tree-based models used fixed
hyperparameters (Table~\ref{tbl:hyperparams}). This is not a fully nested
algorithm comparison, and the search budget may favour the deep-learning
estimates \cite{varma2006bias}. The same caveat applies to the claim that
the random forest is the best model, because that model was also selected
from the observed benchmark. The reported differences should therefore be
read as relative rankings under a fixed evaluation budget, not as unbiased
population estimates of generalisation error.

\subsection{Model-attribution patterns and possible physical
interpretation}\label{grinding-process-insights}

Beyond predictive accuracy, the attribution analyses identify measured
frequency regions associated with the fitted roughness predictions. They do
not establish how surface roughness is physically generated.

\subsubsection{Frequency signatures of surface
generation}\label{frequency-signatures}

The explanation analysis in Section~\ref{sec:feature-attribution} shows a
broad separation between higher-frequency AE and structural-frequency
vibration information, but the dominant bins differ substantially by model
and representation. AE importance is concentrated mainly above 800~kHz, with a substantial
fraction above the historically reported nominal 1~MHz response limit. This pattern is
compatible with high-frequency AE activity associated with grain--workpiece
contacts and micro-fracture events \cite{tonshoff1992acoustic}, although
the component above 1~MHz is not calibrated by an archived sensor response.
These frequencies may encode aspects of grain--workpiece interaction, but
this interpretation requires sensor-response and process-force validation.
For the reproducible current random forest, 34.21\% of out-of-fold AE attribution mass lies
above 1~MHz, outside the historically reported nominal response band; this
component is therefore treated as a model attribution rather than strong
physical evidence.

Vibration attribution is model-dependent: LightGBM TreeSHAP peaks near
20.9~kHz, ResNetVibCNN Grad-CAM peaks near 4~kHz with additional 10--12~kHz
regions, and the current RF out-of-fold TreeSHAP analysis assigns its
strongest vibration importance mainly between approximately 16 and 24~kHz.
Lower-frequency regions receive relatively little
attribution in the evaluated fitted models. These associations are not a
direct agreement between explainers or a causal frequency taxonomy. This
range lies close to the 25.6~kHz Nyquist limit. Because the accelerometer
response, analogue anti-alias filter, and mounting transfer function were
not archived, the attribution and ablation sensitivity cannot distinguish
mechanical content from upper-band acquisition-chain behaviour.
Sensor-transfer-function and fold-stability analyses would be required to
rule out recording-chain artefacts; because no-cut spectra were not
archived, attribution-based physical interpretation remains plausible rather
than experimentally verified.

\subsubsection{Condition 7 as a case study in outlier-driven
failure}\label{condition-seven-case}

Condition~7 is the largest outlier across the evaluated model--configuration
inventory: its mean MAE of 0.234~\(\upmu\)m accounts for 35.2\% of the
summed condition-level mean MAE and is more than
twice that of the next hardest condition (Section~\ref{sec:hard-conditions}).
The failure is not model-specific; every architecture performs worse on
Condition~7, which points to a robust prediction failure rather than a
weakness of any one learner. A critical contributing factor for tree-based
models is the \emph{extrapolation regime}: Condition~7 contains the
highest observed pass-level $R_a$ values (up to 0.381~\(\upmu\)m), which
are absent from the training set when this condition is held out
(Supplementary Table~\ref{tab:supp-logo-extrapolation}). The evaluated constant-leaf random forest cannot produce predictions beyond the range represented by its terminal-leaf training responses and therefore underpredicts near the training maximum of
0.259~\(\upmu\)m (Condition~8, the next-roughest condition), explaining the large
held-out-condition errors. This is a limitation of the evaluated
constant-leaf tree ensembles; it cannot be solved by spectral descriptors
alone unless the training data or model class supports the target range.
Because the design is fractional, this
interpretation is partly confounded with the workpiece-speed/depth-of-cut
pairing and cannot be cleanly isolated from interaction effects.

The training-target-support distance is a post-hoc diagnostic because it
requires the measured test \(R_a\), which is unavailable at prediction
time. It therefore cannot serve as an online warning signal. Deployable
out-of-distribution monitoring must instead use input-space evidence such
as distance in process-parameter or representation space, density/leverage,
ensemble disagreement, or sensor-feature support.
As a descriptive check, the nearest remaining-condition centroid in the
training-fold-standardised AE-dB-z + Vib-dB-z space has a Spearman
association of \(\rho = 0.512\) with the canonical RF LOGO MAE across the 16
conditions (unadjusted \(p=0.043\); leave-one-condition-out range
0.414--0.625; Supplementary Table~\ref{tab:supp-input-ood}). The analogous
nearest process-parameter distance has \(\rho = -0.256\). Thus, the archived
representation contains a potentially useful novelty signal that is not
available from the three programmed parameters alone. This 16-condition,
post-hoc association is not a calibrated online alarm; a deployable method
must fix a threshold without using the unknown test target and be validated
on independent campaigns.

Table~\ref{tbl:condition-seven} compares the process parameters of
Condition~7 with the median of the remaining 15 conditions. Condition~7
combines the highest workpiece speed (80~r/min), a below-median wheel speed
(30~m/s), and the shallowest grinding depth (20~\(\upmu\)m). Condition~7 is
mechanistically ambiguous. Its high workpiece speed and low wheel speed
increase the kinematic chip-thickness proxy, while the shallow depth and
possible wheel-state effects could still produce rubbing, ploughing, or
thermal/wheel-loading behaviour. The current data are insufficient to
distinguish these mechanisms. The kinematic undeformed chip thickness scales as
\begin{equation}
h_m \propto \left(\frac{v_w}{v_s}\right)^{1/2}\left(\frac{a_e}{D_e}\right)^{1/4},
\label{eq:chip-thickness}
\end{equation}
where \(v_w\) and \(v_s\) are expressed in consistent linear surface speed units (using the workpiece outer diameter \(d_w = 214\)~mm to convert \(v_w\) from r/min to m/s), \(a_e\) is the effective depth of cut, and \(D_e\) is the
equivalent wheel diameter \cite{malkin2008grinding, tonshoff1992acoustic}. For this descriptive calculation, \(a_e=a_p\); \(D_e\) and the r/min-to-surface-speed conversion cancel in the ratio because the same wheel and workpiece diameter apply to both cases. Substituting Condition~7 and the component-wise medians gives
\(\sqrt{(80/30)/(40/35)}(20/40)^{1/4}=1.28\).
This kinematic interpretation remains
post-hoc and is not definitively settled. Because our experimental setup
did not archive direct grinding force, spindle power, contact temperature,
wheel loading/dressing state, coolant state, surface images, or
high-resolution spectral difference plots for Condition~7, the exact
material-removal mechanism remains speculative. Future work must integrate
these physical diagnostics to validate regime transitions.

\begin{table}[htbp]
\centering
\caption{Process parameters for Condition~7 versus the median of the
other 15 conditions.}\label{tbl:condition-seven}
\begin{tabular}{@{}lcc@{}}
\toprule
Parameter & Condition~7 & Median of others \\
\midrule
Wheel speed (m/s) & 30 & 35 \\
Workpiece speed (r/min) & 80 & 40 \\
Grinding depth (\(\upmu\)m) & 20 & 40 \\
\bottomrule
\end{tabular}
\end{table}

Condition~10 provides a near-matched diagnostic contrast: it shares the
same workpiece speed and depth of cut as Condition~7 but uses 35 rather
than 30~m/s wheel speed. Its mean measured \(R_a\) is
0.0783~\(\upmu\)m, versus 0.3540~\(\upmu\)m for Condition~7. Figure~\ref{fig:condition7-condition10}
compares the archived mean dB spectra and shows why this target contrast
should be treated as a priority data and machine-state audit, not as proof
of a wheel-speed-only regime transition. Machine logs, dressing/loading
state, coolant state, surface images, force, power, and temperature data
are unavailable for distinguishing the alternatives.
The quantitative cache comparison is consistent with this distinction:
Condition~7 is 2.77~dB higher in the archived 0.1--1.0~MHz AE band and
3.14~dB lower in the archived 2--15~kHz vibration band than Condition~10;
the respective nonparametric pass-bootstrap 95\% intervals are
[1.37, 4.16] and [-4.24, -2.09]~dB (Supplementary
Table~\ref{tab:supp-condition7-condition10-bands}). These are
cache-level contrasts after the recorded signal chain, not calibrated
spectral measurements or evidence for a material-removal mechanism. The
bootstrap resamples repeated passes within each single condition; its
interval quantifies within-condition pass variability and cannot establish
generalisability of the wheel-speed contrast across independently repeated
conditions.

\begin{figure}
\centering
\includegraphics[width=\linewidth]{condition7_condition10_spectra.png}
\caption{Archived-cache spectral comparison for Conditions~7 and 10.
The conditions share workpiece speed and depth of cut but differ in wheel
speed and measured roughness. These descriptive spectra do not identify a
causal mechanism.}\label{fig:condition7-condition10}
\end{figure}

Condition~7 should be treated as a critical failure case for model
robustness and may require either explicit regime identification or
condition-aware uncertainty estimation. In practice it warrants
inspection of wheel wear, dressing state, coolant flow, and possible
thermal damage. It should remain part of the main evaluation unless it is
proven to be invalid data, because difficult or outlier conditions are
precisely the cases that matter most in process monitoring.

\subsection{Reliability and decision-making under
uncertainty}\label{uncertainty-and-decisions}

\subsubsection{Coverage failure of MC-dropout epistemic-only intervals}\label{why-mc-dropout-fails}

MC-dropout produced epistemic-only bands with 50.2\% descriptive target coverage
under full 16-fold LOGO evaluation (319 held-out samples). The intervals
are therefore substantially under-covering. However, predicted standard
deviation is correlated with absolute error at pass level (Pearson
$r=0.488$, Spearman $\rho=0.594$; AUROC for detecting the top 20\%
highest-error samples is 0.818), so MC-dropout has some discrimination
ability but does not provide calibrated predictive intervals. This is not a
formal calibration test of the epistemic posterior variance because the band
omits residual and aleatoric variability. The corresponding 16-condition mean
correlation is Pearson $r=0.492$ (Spearman $\rho=0.732$), but the
condition-cluster bootstrap intervals are wide and the pass-level metrics
must remain descriptive.
Condition~7 has 0\% coverage, showing that it fails on the dominant hard
condition. The earlier 59-sample subset (Conditions 1, 2, and 6)
reported 100\% coverage with a wider mean width (0.073~\(\upmu\)m versus
0.0479~\(\upmu\)m under full LOGO), demonstrating that the apparent quality of
MC-dropout is highly sensitive to which conditions are included. Under the full LOGO evaluation most uncertainty bins are far
below the nominal 95\% level, and the mean interval width is much smaller
than the mean absolute error in the highest-error bins. The failure is
therefore not merely an artefact of the earlier 59-sample subset. It represents a substantial calibration failure for the evaluated CNN, dropout configuration, and grouped-shift setting: dropout variance partially ranks error magnitude but does not
provide calibrated prediction intervals and should not be used alone as a
process-warning signal \cite{ovadia2019can}.

Contributing reasons for the poor predictive coverage include omission of
residual or aleatoric variability, systematic model bias under condition
shift, model misspecification, and limited sensitivity of the evaluated
dropout approximation to out-of-support inputs. If an outlier similar to Condition~7 were presented to the network, the
dropout masks might still produce similar hidden representations, so the
interval could remain narrow even though the prediction would likely be
wrong.

This distinction has direct operational consequences. A plant engineer
relying on MC-dropout intervals would not receive a reliable warning that
Condition~7 is about to produce a large residual. Aggregate coverage is
therefore not sufficient for process control; the calibration of
uncertainty at the sample level matters at least as much as the overall
coverage rate (Figure~\ref{fig:reliability}).

\subsubsection{Practical alternatives for shop-floor
uncertainty}\label{shop-floor-uncertainty}

The failure of both MC-dropout and OOB-scaled RF intervals on Condition~7
indicates that future uncertainty methods must calibrate at condition
level and include frontier conditions. Ensembles of
structurally different models (tree, CNN, scattering) can provide
uncertainty estimates that reflect genuine model disagreement rather than
internal Monte-Carlo noise \cite{lakshminarayanan2017simple}. Conformal
prediction is a candidate approach, but with only 16 conditions the calibration set would be small, and coverage would depend strongly on whether calibration conditions span the same process frontier as Condition~7 \cite{angelopoulos2021gentle}.
Exchangeability can be questioned when conditions differ as strongly as
Condition~7 does from the rest, so adaptive or condition-aware conformal
variants should be considered in practice. For multivariate industrial
targets, copula-based conformal methods provide a way to build joint
prediction regions that respect output correlations
\cite{zhang2023improved}. Evidential regression models, such as Normal-Inverse-Gamma
formulations, or explicit aleatoric/epistemic decomposition could also
separate process noise from model ignorance.
Evaluating these alternatives under the same LOGO protocol is an important
direction for future work.

\subsection{Industrial deployment
implications}\label{industrial-deployment}

The present evidence is retrospective and pass-level: every cached input is
the arithmetic aggregate of 2,910 local maps spanning its recorded pass.
Accordingly, the reported feature-precomputed inference times do not imply
within-pass prediction, a closed-loop update interval, or first-decision
latency. A deployable streaming system would need to fix an aggregation
horizon and hop, then measure raw-signal-to-decision timing and accuracy on
partial or overlapping pass windows.

\subsubsection{Matching model choice to plant
constraints}\label{model-choice-by-constraints}

The deployment trade-offs in Section~\ref{sec:latency-accuracy} and
Section~\ref{sec:footprint-maintainability} can be summarised as a simple
decision rule (Table~\ref{tbl:deployment-choice}). Under the present post-selection benchmark, if the overriding priority
is point-estimate accuracy and the application can tolerate a
pooled median of approximately 8~ms (p95 13.7~ms) for feature-precomputed
model inference, the random
forest on fused AE and vibration dB-z descriptors is one leading
point-estimate candidate alongside the nearly tied log-mel RF by mean MAE.
The current nested dB-z result has the lower maximum fold error
(0.1402 versus 0.2061~\(\upmu\)m), so dB-z is the better observed tail-risk
candidate in this dataset (its
superiority is not statistically conclusive versus most deep-learning
alternatives under the primary test family). It can be audited through the
out-of-fold RF TreeSHAP analysis reported for the same dB-z pipeline. If the application
requires sub-5~ms median inference on an edge device or must run at high
cached-feature model-call rates, ResNetVibCNN provides the lowest median and
p95 latency, while ResNetAECNN is the AE-only CNN alternative. LightGBM and
bilinear fusion also have pooled medians near or below 5~ms but different
accuracy and interpretability trade-offs. Both CNN choices have MAEs
approximately 0.0068--0.0100~\(\upmu\)m above the historical best RF point
estimate; most corresponding pairwise differences are not statistically
conclusive after multiplicity correction. Their decisions are also less
directly auditable under the explanation methods evaluated here.

\begin{table}[htbp]
\centering
\caption{Recommended model--configuration pair by deployment
priority. All latency figures are \emph{feature-precomputed model-inference
only}; true online latency must additionally account for the full-pass temporal aggregation used by the current cache, signal acquisition, STFT/dB-z/log-mel preprocessing, feature
scaling, and communication/edge-hardware overhead.}
\label{tbl:deployment-choice}
\begin{tabular}{@{}lll@{}}
\toprule
Priority & Recommended model & Configuration \\
\midrule
Lowest observed point-estimate group & Random forest & \texttt{AE-dB-z + Vib-dB-z} or \texttt{AE-log-mel + Vib-log-mel} \\
Better observed nested worst-fold robustness & Random forest & \texttt{AE-dB-z + Vib-dB-z} \\
Lowest median and p95 model-only latency & ResNetVibCNN & \texttt{Vib-dB} \\
Out-of-fold attribution for recommended model & Random forest & \texttt{AE-dB-z + Vib-dB-z} \\
Exact local TreeSHAP where individual examples are required & LightGBM & \texttt{AE-dB + Vib-dB} \\
Observed low-error baseline under limited data & Random forest & \texttt{AE-dB-z + Vib-dB-z} \\
\bottomrule
\end{tabular}
\end{table}

\noindent\textbf{Latency caveat.} The reported inference latencies are
\emph{feature-precomputed model-inference times only} and do not represent
end-to-end online latency. The complete acquisition-to-decision pipeline would
additionally require a chosen partial-pass aggregation horizon and hop,
signal acquisition, signal preprocessing (STFT, dB-z/log-mel computation,
feature scaling), network communication, and edge-hardware boot-up. The
current cached descriptor instead aggregates the full recorded pass
(approximately 14.6~s in the raw AE record). Industrial claims
about real-time capability must specify the full pipeline budget; the present
latency figures serve only as a relative ranking of model-inference overhead.

LightGBM offers a middle ground in accuracy and its p95 latency is
5.0~ms in the fold-complete controlled rerun, which keeps it within the same
order-of-magnitude latency band as the RF configurations while remaining
less accurate than the best RF. Bilinear fusion is not competitive: it is the largest deep-learning entry without any
point-estimate accuracy advantage over the simpler CNNs.

\subsubsection{Human-in-the-loop, interpretability, and
retraining}\label{human-in-the-loop}

Auditable models may be easier to inspect because their frequency-bin attributions and split rules can be examined, but this expected advantage was not tested through a user study, retraining-time benchmark, or plant deployment. TreeSHAP attributions let engineers inspect whether a
prediction is associated with physically plausible frequency regions
rather than with sensor drift or anomalous wheel loading, although this
audit is post-hoc and does not prove causal reliance. When a new wheel
grade, workpiece material, or dressing strategy is introduced, the random
forest is expected to require less tuning than a deep CNN, but this
maintainability claim was not measured here. For production environments
where models must be updated frequently, this expected advantage should be
tested alongside point-estimate accuracy.

\subsection{Limitations and future
work}\label{limitations-and-future-work}

\subsubsection{Dataset scale and
generalisability}\label{scale-and-generalisability}

The study is based on 319 samples from 16 grinding conditions. While this
collection is large enough to expose clear performance differences and
condition-level failure modes, it is not large enough to guarantee
generalisation to every wheel grade, workpiece material, dressing protocol,
or coolant formulation. Future work should expand the dataset across these
dimensions and report external validation on independently collected data.

\subsubsection{Statistical power and inference limitations}\label{power-limitations}

With only 16 condition-level folds, most pairwise model comparisons are underpowered to detect subtle performance differences under Holm--Bonferroni correction. The experimental design can detect only relatively large paired condition-level effects; non-significance is not evidence of equivalence. Consequently, while the random forest's point-estimate advantage is suggestive, the lack of statistically significant differences against many deep learning baselines must be interpreted with caution.

\subsubsection{Reproducibility, preprocessing, and measurement
caveats}\label{reproducibility-caveats}

Several reproducibility caveats bound the scientific claims. Only the
final averaged $R_a$ value per pass is archived; the raw profilometer
traces and trace-level repeatability were not retained, so gauge R\&R
	and experimental uncertainty cannot be computed. The best reported MAE
	(0.020~\(\upmu\)m) is only about twice the manufacturer-quoted
	repeatability (0.01~\(\upmu\)m) and is below the expanded uncertainty
	scale of 0.025~\(\upmu\)m. It should therefore be interpreted as error
	against the averaged measured \(R_a\), not error against latent true
	roughness. Because trace-level repeatability is unavailable, the latent
	error relative to latent roughness cannot be separated from
	measurement uncertainty. Differences among the top RF
variants are smaller than the available measurement-resolution evidence,
so their ranking is exploratory. No-cut background spectra were
not archived, so frequency attributions cannot be verified against
background-noise measurements and remain plausible rather than
experimentally confirmed. A strict per-sample/channel normalisation ablation for the top RF
configuration changes the overall LOGO MAE only slightly (0.0210
versus 0.0213~\(\upmu\)m), but the worst-condition MAE increases without
per-sample normalisation (Supplementary
Table~\ref{tab:supp-psnorm-ablation}), so the ablation supports
mean-level robustness but not necessarily tail robustness. The top RF
result does not therefore appear to depend on removing absolute
AE/vibration energy for mean MAE. Physics-informed features are computed by
the implemented algorithms and parameters documented in Supplementary
Table~\ref{tab:physics_features}; any cached intermediate-derived
preprocessing must be preserved for full reproducibility. Historical campaign
notes report a nominal 100~kHz--1~MHz AE response band, but the sensor
specification and transfer-function record are unavailable; AE
attributions above 1~MHz (34.21\% of the current RF's out-of-fold AE importance)
remain sensor-response-limited and must not be treated as validated grinding physics. In this high-frequency regime, the measured content is heavily shaped by sensor-mounting, couplant resonance, and the sensor's own transfer function rather than raw workpiece interactions. Because our experimental setup did not characterize the sensor transfer function, collect no-cut spectra, or measure mount/couplant resonance, we cannot physically validate the attributions above 1~MHz. The nominal 100~kHz--1~MHz band is a cautious reporting convention, not a calibrated physical-interpretation range; high-frequency features serve solely as empirical statistical predictors.

\subsubsection{Wheel-state dynamics and cyclostationary
features}\label{wheel-state-dynamics}

The current feature set captures frequency content but does not explicitly
model wheel-pass modulation or cyclostationary dynamics. As a wheel wears
or loads, the periodicity of grain engagements changes, and
cyclostationary features---cyclic spectral correlation, averaged cyclic
periodograms, or spectral kurtosis---could track this evolution more
directly than static spectrograms \cite{antoni2007cyclic, antoni2004cyclostationary}.
Including such features is a natural next step, particularly for
Condition~7-style outliers where wheel state may be the dominant factor.

\subsubsection{Robustness of feature and validation
choices}\label{feature-validation-robustness}

The cross-validation-scheme comparison shows that random sample-level
splits give MAEs of roughly 0.005~\(\upmu\)m, about four times lower than
grouped LOGO. Because this value is below the manufacturer-quoted
profilometer repeatability and primarily reflects within-condition similarity, it
should not be interpreted as practically meaningful performance. This
confirms that condition-level grouping is essential and that the modest
differences among top models are only meaningful under grouped evaluation.
The current descriptors are temporal means over 2,910 archived local windows per pass; they do not retain early/middle/late temporal ordering. The retrospective early/middle/late-third sensitivity and ordered three-window concatenation have similar mean MAEs, but fixed-distance and streaming-dispersion features remain necessary to test whether transients add information beyond the pass-level mean. The sensor-modality ablation shows that
vibration alone achieves 0.0217~\(\upmu\)m, nearly matching the fused
0.0210~\(\upmu\)m result, while AE alone degrades to 0.0282~\(\upmu\)m.
This comparison uses the exact 14-condition current benchmark split. AE therefore
provides a secondary, condition-specific benefit rather than a dominant
signal. The frequency-band ablation shows that removing AE content above
1~MHz or the vibration 2--15~kHz chatter band individually changes MAE by
less than 0.0001~\(\upmu\)m; only removing both bands together raises MAE
to 0.0208~\(\upmu\)m. This supports treating the >1~MHz AE attribution as
a model response pattern rather than calibrated physical energy: the
band is not required for near-optimal accuracy, even though the RF
assigns it substantial importance. The complementary retrained analysis
finds no interval-supported condition-level change after removing AE >1~MHz,
vibration 2--15~kHz, or the narrower vibration 16--24~kHz region. In
contrast, removing the complete vibration region above 15~kHz raises MAE
from 0.0210 to 0.0285~\(\upmu\)m, with a descriptive paired 95\% interval
that excludes zero. The fitted RF therefore uses predictive information
distributed across the high-frequency vibration region; this result does
not establish a causal grinding mechanism. A stricter dB-source analysis
deletes the bands before recomputing the retained-bin z-score and again finds
an increase for removal of >15~kHz (0.0209 to 0.0255~\(\upmu\)m). Therefore,
the finding is not explained solely by indirect encoding through the original
full-spectrum normalisation.

\subsubsection{Broader uncertainty-quantification
methods}\label{broader-uq}

The MC-dropout analysis is limited to one CNN architecture and shows poor
per-sample calibration. The full-LOGO MC-dropout values reported in
Section~\ref{uncertainty-quantification} supersede the earlier 59-sample
subset, which is retained only to show sensitivity to subset choice. We
also tested four interval strategies for a separate 15-condition RF uncertainty variant (global OOB scaling, condition-aware k-NN calibration,
split conformal with one held-out calibration condition, and a naive
residual interval). None covers Condition~7, indicating that the largest
outlier lies outside the support of in-distribution residuals. Future
work should compare conformal prediction
\cite{angelopoulos2021gentle, zhang2023improved}, deep ensembles
\cite{lakshminarayanan2017simple}, and evidential or explicitly decomposed
uncertainty estimates under the same LOGO protocol, with explicit
calibration reserves that span the parameter-space frontier. The goal is
an uncertainty estimate that not only covers the aggregate error but also
flags condition-specific outliers before they enter the control loop.

\subsubsection{Benchmarking and deployment caveats}

The model ranking is the outcome of a non-nested/post-selection benchmark:
the best random forest and the best deep-learning configurations were
selected from the same observed results, and the deep-learning entries
were further tuned by random search on the outer LOGO folds. The reported
differences should be read as relative rankings under a fixed evaluation
budget, not as unbiased population estimates. This caveat applies to the
selected best RF as well as to the deep-learning entries. A current-cache,
inner-grouped rerun gives 0.0207~\(\upmu\)m for dB-z and 0.0206~\(\upmu\)m
for log-mel RF descriptors; an earlier nested log-mel artifact gives
0.0455~\(\upmu\)m and is not reproduced by the current pipeline
(Supplementary Table~\ref{tab:supp-nested-logo}). This provenance discrepancy
reinforces that dB-z and log-mel should be treated as a leading observed group,
not a uniquely ordered pair. Random 5-fold CV and
random 80/20 splits give MAEs of roughly 0.0053~\(\upmu\)m, so validation
schemes that allow condition leakage between training and test sets are
overly optimistic for this dataset; grouped LOGO or parameter-group LOGO
is required for an honest estimate of generalisation. Four RF interval
calibration strategies (global OOB, condition-aware k-NN, split
conformal, naive residual) all fail to cover Condition~7, confirming that
the dominant outlier is not captured by in-distribution residual
patterns. The latency analysis is limited to feature-precomputed
model-inference timing on a single CPU core with the `powersave'
governor; it excludes signal acquisition, the full-pass temporal aggregation used by the current cache, STFT/dB-z/log-mel computation, feature scaling, and communication
	overhead. A full end-to-edge deployment assessment remains future work, and
	the reported timings do not imply full edge-deployment readiness. Robust
	latency benchmarking should repeat the measurements in multiple blocks
	under a fixed CPU performance governor and report p90, p95, p99, and
	maximum latency before using any model for hard real-time conclusions.
